\chapter{Context}

    \section{Scope of Project}

        \paragraph{} This thesis investigates the design and implementation of an ontology-driven, knowledge-graph-enabled systems engineering framework that integrates structured semantic models with the contextual reasoning capabilities of large language models (LLMs). Anchored within Sunswift Racing, a world-leading student engineering organisation focused on the research, development, and demonstration of next-generation solar-electric vehicles, the project aims to address the increasing complexity of regulatory compliance and design traceability across modern automotive systems.

        \vspace{1cm}
        \begin{figure}[h]
            \centering
            \includegraphics[width=1\linewidth]{Photos/SR8.jpeg}
            \caption{SR8 Rendering}
            \label{fig:rendering}
        \end{figure}

        \newpage
        \paragraph{} Sunswift Racing has, since 1996, produced seven generations of experimental solar-electric vehicles, achieving international recognition through the Bridgestone World Solar Challenge (BWSC), FIA World Records, and Guinness World Records. The upcoming Sunswift 8 (SR8) platform represents a paradigm shift: the project is intended to be the world’s first road-legal, student-built, tri-brid electric vehicle. This ambition demands unprecedented levels of systems engineering rigour, cross-disciplinary coordination, and regulatory alignment.

        \paragraph{} The scale of SR8’s engineering challenge is substantial: over 80 team members, distributed across \textbf{16 technical departments}, working on \textbf{189 concurrent projects}, interfacing through more than \textbf{5,000 sub-system interactions}, and collectively producing and maintaining approximately \textbf{600,000 requirements}, \textbf{100,000+ components}, \textbf{1.5 TB of simulation data}, and nearly \textbf{1 TB of engineering drawings}. Much of this knowledge is dispersed across heterogeneous formats, spreadsheets, reports, test plans, engineering notes, simulations, and regulatory documents, resulting in fragmentation, duplication, and limited traceability.

        \vspace{1cm}
        \begin{table}[h]
            \centering
            \renewcommand{\arraystretch}{1.2}
            \begin{tabularx}{\textwidth}{|c|X|X|}
            \hline
            \textbf{Quantity} & \textbf{Description} & \textbf{Department} \\
            \hline
            \textbf{80} & Number of Members & Whole Team \\ \hline
            \textbf{16} & Number of Departments \newline (Technical + Operational) & Whole Team \\ \hline
            \textbf{189} & Number of Projects & Whole Team \\ \hline
            \textbf{5,000+} & Number of Interactions \newline (Interfaces) & Whole Team \\ \hline
            \textbf{100,000+} & Number of Components \newline (Procurement) & Whole Team \\ \hline
            \textbf{113} & Number of Australian Design Rules & Whole Team \\ \hline
            \textbf{600,000+} & Number of Requirements \newline (Internal + External) & Whole Team \\ \hline
            \textbf{4 TB} & Amount of Documents, Reports, \newline Excel Sheets & Whole Team \\ \hline
            \textbf{1.5 TB} & Mechanical Simulations & INT, C\&B, PWT, VD, ID \\ \hline
            \textbf{1 TB} & Electrical Drawings & EMB, ENS, ALT \\ \hline
            \textbf{500 GB} & High-Level \& Low-Level Code & EMB, TSS, SYS \\
            \hline
            \end{tabularx}
            \caption{SR8 Complexity}
        \end{table}

        \paragraph{} Within this context, the thesis situates itself as a response to a broader systems-engineering need: to transform disconnected documentation into a unified, queryable, and explainable knowledge infrastructure. The project focuses particularly on the Australian Design Rules (ADR's), the national regulatory framework governing vehicle safety, performance, construction, and compliance. By formalising ADR's into an ontology-based knowledge graph and coupling this structure with LLM-assisted reasoning, the research aims to enable automated requirement extraction, semantic alignment, risk propagation, and model-based compliance verification across the Sunswift design eco-system.

    \newpage
    \section{Requirements}

        \paragraph{} The requirements driving this research arise from the intersection of regulatory complexity, organisational scale, and the increasing demand for model-based engineering workflows. Specifically, the framework developed in this thesis must:

        \begin{enumerate}
          \item \textbf{Represent Regulatory Knowledge Formally:} \\ The Australian Design Rules must be modelled as a machine-interpretable ontology, capturing requirements, constraints, justifications, exemptions, applicability conditions, and cross-references. This includes the ability to express hierarchical rule structure, sub-system interactions, compliance dependencies, and mapping across Sunswift’s projects.
          \item \textbf{Enable Semantic Traceability At Scale:} \\ Given the magnitude of Sunswift’s multi-department engineering operations, the system must support robust traceability across projects, components, interfaces, tests, and documentation. Requirements must be linkable to design decisions, verification activities, and sub-system responsibilities.
          \item \textbf{Integrate Structured \& Un-Structured Data Sources:} \\ A central requirement is the ability to unify heterogeneous engineering information, PDF reports, test logs, simulation outputs, spreadsheets, subsystem documents, risk assessments, and ADR text, into a coherent knowledge representation. This necessitates automated information extraction, natural-language understanding, and schema alignment.
          \item \textbf{Provide Explainable, Context-Aware Reasoning:} \\ Engineers require not only compliance verification, but also explanations of why a requirement applies, how it propagates through the system, and what dependencies are affected. LLM-driven reasoning must complement the deterministic structure of the ontology to support nuanced interpretation and expert-level contextualisation.
          \item \textbf{Support Automated Risk Identification \& Decision Guidance:} \\ The system should incorporate risk modelling, enabling the detection of design gaps, non-compliance risks, and potential system-level impacts. A hybrid recommender system, combining LLM inference, knowledge-graph reasoning, and reinforcement learning, must be capable of suggesting improvements, alternative design paths, or additional verification steps.
          \item \textbf{Enable Visual, Interactive Exploration For Engineers:} \\ A key usability requirement is the development of an interactive interface enabling engineers to navigate regulatory knowledge as interconnected nodes rather than static text. This aligns with human-computer interaction principles and modern digital-thread practices in model-based systems engineering.
        \end{enumerate}
        
        \paragraph{} Collectively, these requirements define the functional and non-functional expectations for a unified regulatory-aligned systems engineering framework capable of supporting high-complexity automotive development.

    \section{Engineering Problem}

        \paragraph{} Despite Sunswift’s longstanding engineering excellence and strong culture of innovation, the current management of regulatory knowledge, particularly the Australian Design Rules, remains highly fragmented. Although the team produces immense volumes of structured and unstructured design information, this knowledge is distributed across disconnected documents, local storage systems, and department-specific formats without a unified semantic or systems-engineering structure.
        
        \paragraph{} Each ADR specifies detailed requirements, constraints, and justifications that influence sub-system design decisions across multiple departments. However, these rules exist primarily as large, manually reviewed PDF documents that provide no inherent traceability, contextual metadata, or project-level mappings. Engineers must search through lengthy text documents to determine applicability, dependencies, or compliance pathways. This manual process introduces several systemic challenges:

        \begin{enumerate}
          \item \textbf{Missed Dependencies:} \\ Requirements that apply to multiple sub-systems or projects are often overlooked due to inconsistent naming, formatting, or document organisation.
          \item \textbf{Redundant Or Conflicting Interpretations:} \\ Departments independently transcribe or reinterpret ADR's, leading to duplication and semantic drift.
          \item \textbf{Stagnant Or Out-Dated Documentation:} \\ Updates to requirements rarely propagate throughout the organisation, causing inconsistencies in design decisions, verification plans, and audit documentation.
          \item \textbf{Inefficient Compliance Workflows:} \\ Verification, testing, certification, and review processes rely heavily on human interpretation of unstructured text, increasing risk and engineering overhead.
          \item \textbf{Lack of Model-Based Reasoning:} \\ Without a semantic representation of requirements, the organisation cannot support data-driven traceability, automated reasoning, or predictive compliance analysis.
        \end{enumerate}

        \paragraph{} This fragmented state is exacerbated by the scale of SR8: with hundreds of thousands of requirements, tens of thousands of components, and terabytes of engineering artefacts, the absence of a unified knowledge structure hinders Sunswift’s ability to maintain consistency, ensure safety, and integrate regulatory compliance throughout the design-to-validation pipeline.

        \newpage
        \paragraph{} The engineering problem can be summarised along the five classical interrogatives:

        \begin{itemize}
            \item \textit{What:} A pervasive fragmentation of regulatory knowledge, particularly the ADR's, across documents, departments, and design artefacts.
            \item \textit{Why:} The absence of a standardised ontology or traceability framework leads to inconsistent interpretation, limited searchability, and manual compliance verification.
            \item \textit{Where:} Across 189 projects and 16 departments, where requirements appear in multiple ADR's, often worded differently, duplicated, or detached from the systems they constrain.
            \item \textit{When:} Throughout the entire engineering life-cycle, from conceptual design and sub-system integration to testing, verification, and audit, where compliance decisions must be continuously updated and validated.
            \item \textit{How:} Currently through manual review of PDF's, spreadsheets, and reports, without automated extraction, semantic mapping, or system-level reasoning, resulting in inefficiency, errors, and incomplete traceability.
        \end{itemize}

        \paragraph{} Addressing this problem requires a shift from document-centric engineering to a knowledge-centric, model-based paradigm. By formalising regulatory text into an ontology-driven knowledge graph and augmenting this structure with large language models for contextual reasoning, this thesis proposes an integrated framework capable of enabling adaptive compliance management, semantic traceability, risk propagation, and explainable decision support. In doing so, this thesis seeks to transform Sunswift’s design eco-system into a scalable, intelligent, and future-proof systems engineering environment.
